\documentclass[11pt]{article}
\usepackage[a4paper,margin=2cm]{geometry}
\usepackage{amsmath,amssymb}
\usepackage[hidelinks]{hyperref}
\allowdisplaybreaks
\title{Qpu.General.Iso}
\author{Tsvetan Rouschev}
\date{@uuidna/qpu v1.1.1, CC-BY-NC-ND-4.0}
\begin{document}
\maketitle
\noindent\url{https://github.com/uuidna/qpu/blob/v1.1.1/lean/Qpu/General/Iso.lean}

\paragraph{\texttt{Iso.calendar\_cycle}}
ISO 8601: four hundred years hold ninety-seven leap years and a whole number of weeks.
\[ \mathrm{Iso.leap}\ 2024 = \mathrm{true} \land \mathrm{Iso.leap}\ 1900 = \mathrm{false} \land \mathrm{Iso.leap}\ 2000 = \mathrm{true} \land (\mathrm{List.filter}\ \mathrm{Iso.leap} (\mathrm{List.range}\ 400))\mathrm{.length} = 97 \land (400 \cdot 365 + 97) \bmod 7 = 0 \]
\noindent axioms: none; proof: \url{https://github.com/uuidna/qpu/blob/v1.1.1/lean/Qpu/General/Iso.lean\#L178}

\paragraph{\texttt{Iso.clauses\_count}}
The register holds eighty-six terms: the sum of its groups.
\[ \mathrm{Iso.clauses.length} = 86 \land (\mathrm{List.map} (\lambda \ x \mapsto \mathrm{match}\ x\ \mathrm{with} | (\mathrm{fst}, \mathrm{fst\_1}, n) \mapsto n) \mathrm{Iso.groups})\mathrm{.sum} = 86 \]
\noindent axioms: \texttt{propext}; proof: \url{https://github.com/uuidna/qpu/blob/v1.1.1/lean/Qpu/General/Iso.lean\#L170}

\paragraph{\texttt{Iso.clauses\_numbered}}
The numbering is whole: every group runs from one to its last term, in order, with no gap and no repeat.
\[ \mathrm{List.map} (\lambda \ c \mapsto (\mathrm{c.group}, \mathrm{c.index})) \mathrm{Iso.clauses} = \mathrm{List.flatMap} (\lambda \ x \mapsto \mathrm{match}\ x\ \mathrm{with} | (g, \mathrm{fst}, n) \mapsto \mathrm{List.map} (\lambda \ i \mapsto (g, i + 1)) (\mathrm{List.range}\ n)) \mathrm{Iso.groups} \]
\noindent axioms: none; proof: \url{https://github.com/uuidna/qpu/blob/v1.1.1/lean/Qpu/General/Iso.lean\#L165}

\paragraph{\texttt{Iso.country\_codes}}
ISO 3166-1: a country has a code of two letters, a code of three letters and a number of three digits.
\[ \mathrm{Iso.letterCodes}\ 2 = 676 \land \mathrm{Iso.letterCodes}\ 3 = 17576 \land 10 ^{3} = 1000 \]
\noindent axioms: \texttt{propext}; proof: \url{https://github.com/uuidna/qpu/blob/v1.1.1/lean/Qpu/General/Iso.lean\#L192}

\paragraph{\texttt{Iso.currency\_codes}}
ISO 4217: a currency has a code of three letters and a number of three digits.
\[ \mathrm{Iso.letterCodes}\ 3 = 17576 \land 10 ^{3} = 1000 \]
\noindent axioms: \texttt{propext}; proof: \url{https://github.com/uuidna/qpu/blob/v1.1.1/lean/Qpu/General/Iso.lean\#L186}

\paragraph{\texttt{Iso.language\_codes}}
ISO 639-1: a language has a code of two letters.
\[ \mathrm{Iso.letterCodes}\ 2 = 676 \]
\noindent axioms: \texttt{propext}; proof: \url{https://github.com/uuidna/qpu/blob/v1.1.1/lean/Qpu/General/Iso.lean\#L189}

\paragraph{\texttt{Iso.message\_identifier}}
ISO 20022: a message is named by four letters for its business area, three digits for the message, three for the variant and two for the version, as in `pacs.008.001.08`: twelve characters between the dots.
\[ 4 + 3 + 3 + 2 = 12 \land "\mathrm{pacs.008.001.08}"\mathrm{.length} = 12 + 3 \]
\noindent axioms: \texttt{propext}, \texttt{Classical.choice}, \texttt{Quot.sound}; proof: \url{https://github.com/uuidna/qpu/blob/v1.1.1/lean/Qpu/General/Iso.lean\#L205}

\paragraph{\texttt{Iso.quality\_characteristics}}
ISO/IEC 25010:2023: the product quality model has 9 characteristics, one more than the 8 of the edition before.
\[ 8 + 1 = 9 \]
\noindent axioms: none; proof: \url{https://github.com/uuidna/qpu/blob/v1.1.1/lean/Qpu/General/Iso.lean\#L212}

\paragraph{\texttt{Iso.quality\_clauses}}
ISO 9001:2015: of its 10 clauses the first 3 set no requirement, so 7 clauses carry the requirements.
\[ 10 - 3 = 7 \]
\noindent axioms: none; proof: \url{https://github.com/uuidna/qpu/blob/v1.1.1/lean/Qpu/General/Iso.lean\#L201}

\paragraph{\texttt{Iso.security\_controls}}
ISO/IEC 27001:2022: Annex A lists 93 controls in four themes, 37 organizational, 8 people, 14 physical and 34 technological.
\[ 37 + 8 + 14 + 34 = 93 \]
\noindent axioms: none; proof: \url{https://github.com/uuidna/qpu/blob/v1.1.1/lean/Qpu/General/Iso.lean\#L198}

\paragraph{\texttt{Iso.standards\_count}}
The register names twenty standards, and twelve of them have a witness.
\[ \mathrm{Iso.standards.length} = 20 \land (\mathrm{List.filter} (\lambda \ x \mapsto \mathrm{x.witness.isSome}) \mathrm{Iso.standards})\mathrm{.length} = 12 \]
\noindent axioms: none; proof: \url{https://github.com/uuidna/qpu/blob/v1.1.1/lean/Qpu/General/Iso.lean\#L248}

\paragraph{\texttt{Iso.token\_identifier}}
ISO 24165: a digital token identifier is nine characters, eight of the base number and one check character, as in `4H95J0R2X`.
\[ 8 + 1 = 9 \land "\mathrm{4H95J0R2X}"\mathrm{.length} = 9 \]
\noindent axioms: \texttt{propext}, \texttt{Classical.choice}, \texttt{Quot.sound}; proof: \url{https://github.com/uuidna/qpu/blob/v1.1.1/lean/Qpu/General/Iso.lean\#L209}

\paragraph{\texttt{Iso.tuning\_octaves}}
ISO 16: the tuning note is 440 hertz, and an octave doubles or halves it.
\[ 440 \cdot 2 = 880 \land 440 / 2 = 220 \]
\noindent axioms: none; proof: \url{https://github.com/uuidna/qpu/blob/v1.1.1/lean/Qpu/General/Iso.lean\#L195}

\paragraph{\texttt{Iso.vocabulary\_whole}}
ISO/IEC 4879: every term of the register is numbered and witnessed.
\[ \mathrm{Iso.clauses.length} = (\mathrm{List.map} (\lambda \ x \mapsto \mathrm{match}\ x\ \mathrm{with} | (\mathrm{fst}, \mathrm{fst\_1}, n) \mapsto n) \mathrm{Iso.groups})\mathrm{.sum} \]
\noindent axioms: \texttt{propext}; proof: \url{https://github.com/uuidna/qpu/blob/v1.1.1/lean/Qpu/General/Iso.lean\#L215}
\end{document}
